\honest{Savanna Poets}{Eliezer Yudkowsky}{2008}

I open with Feynman's footnote on the stars. Poets say science reduces them to mere gas; nothing is ``mere''; he too sees and feels the stars, and sees more: million-year-old light, a vast pattern of which he is part, his stuff perhaps belched from some forgotten star. ``It does not do harm to the mystery to know a little about it.'' It ends: ``What men are poets who can speak of Jupiter if he were like a man, but if he is an immense spinning sphere of methane and ammonia must be silent?''

I treat that as a real question with a real answer. A Jupiter like us can fall in love, strive and weep. A sphere of gas is harder to make us feel for. Storytellers say the Great Stories are timeless, that Sophocles and Shakespeare could have swapped centuries without a jolt. The strongest emotions are universal: Donald Brown found marriage, incest avoidance, motherly love, sibling rivalry, music, envy, dance, storytelling, healing magic and poetry in all, or nearly all, the cultures studied. ``No one who knows anything about evolutionary psychology could be expected to deny it'': these emotions are engraved in ``brain and DNA.'' \nb{Brown's list shows what every culture has, not why; it also lists weapons and units of time. The post asserts the step to DNA and does not argue it.} You could probably tell \textsc{Hamlet} around a campfire on the savanna.

So Keats, told that the rainbow is light scattered by raindrops, felt a loss. Raindrops don't dance. The Bible's rainbow is God's sign, after a flood that drowned the guilty and their ``horribly guilty babies,'' that he will not do it again, at least not with water. Keats would be shocked to see that story contradicted, especially with no insight to replace it. And even if he knew the mathematics, a tale of bloodthirsty murder and smiling insanity cannot be matched by refraction: raindrops don't scream when they die. What science gives back never matches the drama of the original delusion, because the equations are not about strong emotions. This is ``the strongest rejoinder I can think of'' that a Romantic poet could have made to Feynman, though I cannot remember hearing anyone make it.

My answer: Jupiter need not be like a human, because humans are. Love and hate have not left the universe. The Great Stories happen every day among six billion people: every day someone kills for revenge, someone kills a friend by mistake, and a hundred thousand people fall in love. And you can write fiction about humans, not Jupiter.

Nor are the Great Stories timeless, because the species is not. Go far enough back and no one understands \textsc{Hamlet}; further, and there are no brains. \nb{The poets' claim, as the post reports it, was that Sophocles and Shakespeare could have swapped centuries. For that span the post agrees with them.} I ``most sincerely doubt'' that we have another thousand years in our current form, not in sadness, since I think we can do better. Losing the Great Stories entirely would differ very little from the Sun falling into a black hole, but they have been told over and over, and I would like some of them to change; for me the story called ``Goodbye'' has lost its charm. ``And they lived happily ever after'' seems worth trying once.

The stories should diversify as humankind grows up, and when we find strangeness we should respect it enough to tell its story truly. A poet who writes an ode to the real Jupiter writes something original; writing Jupiter as a human impoverishes our map, forcing it into the stories already told of Earth. James Thomson praised the rainbow for what it really is, in his poem on Newton; whether or not it grips like \textsc{Lamia}, tales of love and loss were already old in Greece, and until the rainbow was understood as something other than human-shaped magic, its true story could not be put into poetry. \nb{Thomson's poem, of 1727, says Newton ``Untwisted all the shining robe of day,'' and describes the colours ``adown the watery bow.'' It came about ninety years before \textsc{Lamia}, which is itself a Greek tale that Keats took from Burton's retelling of Philostratus.}

Real science fiction was once distinguished from space opera as a story that could not be moved to the Old West or the Middle Ages without loss, because the science is part of the plot. Poets who can speak of Jupiter only as a man are ``savanna poets'': they can tell only the stories that made sense around a campfire ten thousand years ago.
